\usepackage{amsmath}
\usepackage{amsthm}
% \usepackage[notref,notcite]{showkeys}
%\usepackage{hyperref}
\usepackage{graphicx}
% \usepackage{slashed}
% \usepackage{amscd}
% \usepackage{tensor}
\usepackage{amssymb}
% \newcommand{\bra}[1]{\ensuremath{\langle #1 |}}
% \newcommand{\ket}[1]{\ensuremath{| #1 \rangle}}

% \usepackage{esint}
\usepackage{enumitem}

% \usepackage{geometry}

% \usepackage{accents}
\usepackage[dvipsnames]{xcolor}
%\usepackage{mathabx}
%\usepackage{amsmath, amscd}
%\usepackage{ngerman}
% \usepackage[utf8]{inputenc}	% Um Umlaute tippen zu können. (Dieses Dokument ist im UTF8 Format angelegt. Alternativen für Mac: latin1, applemac.
%\usepackage{refcheck}
%\usepackage{auto-pst-pdf} % langsame Version, alle Bilder werden als PDF erstellt
% \usepackage[off]{auto-pst-pdf} % schnelle Version, Bilder müssen dazu schon vorliegen
%\usepackage[usenames,dvipsnames]{pstricks}
%\usepackage{epsfig}
% \usepackage{pst-grad} % For gradients
% \usepackage{pst-plot} % For axes
% \usepackage[mathscr]{eucal}
% \textheight 22cm
% \textwidth 14.38cm

% \oddsidemargin=0.9cm
% \evensidemargin=0.9cm
% \topmargin=-0.5cm
% \numberwithin{equation}{section}
% \allowdisplaybreaks[1]

%\newcommand{SetFigFont}[3]{}

%\author[J.\ Kleiner]{Johannes Kleiner \\ \\ January 2018}
%\thanks{Supported in part by the Deutsche Forschungsgemeinschaft.}
%\address{Institut f\"ur Theoretische Physik \\ Leibniz Universit\"at \\ D-30167 Hannover \\ Germany}
%\email{johannes.kleiner@itp.uni-hannover.de}

\theoremstyle{definition}
\newtheorem{Def}{Definition}
\newtheorem{definition}[Def]{Definition}
\newtheorem{Thm}[Def]{Theorem}
\newtheorem{Prp}[Def]{Proposition}
\newtheorem{Lemma}[Def]{Lemma}
\newtheorem{Remark}[Def]{Remark}
\newtheorem{Corollary}[Def]{Corollary}

\newtheorem{Assumption}[Def]{Assumption}
\newtheorem{question}[Def]{Question}

\newtheorem{Example}[Def]{Example}
\newtheorem{example}[Def]{Example}
\newtheorem{examples}[Def]{Examples}


\newcommand{\Decomps}{\mathbb{D}} %All decompositions of object in a monoidal category
\newcommand{\nullexp}{\emptyset}
\newcommand{\Thanks}{\vspace*{.5em} \noindent \thanks}
\newcommand{\intscaling}{\iota}
\newcommand{\QEDrem}{\ \hfill $\Diamond$}
\newcommand{\spc}{\;\;\;\;\;\;\;\;\;\;}
\newcommand{\la}{\langle}
\newcommand{\ra}{\rangle}
%\newcommand{\bra}{\,<\!\!}
%\newcommand{\ket}{\!\!>\,}
% \newcommand{\bra}{\mathopen{<}}
% \newcommand{\ket}{\mathclose{>}}
\newcommand{\Sl}{\mbox{$\prec \!\!$ \nolinebreak}}
\newcommand{\Sr}{\mbox{\nolinebreak $\succ$}}
\newcommand{\ketaux}{\!\!>_\text{aux}\,}
\newcommand{\C}{\mathbb{C}}
\newcommand{\R}{\mathbb{R}}
\newcommand{\1}{\mbox{\rm 1 \hspace{-1.05 em} 1}}
\newcommand{\Z}{\mathbb{Z}}
\newcommand{\N}{\mathbb{N}}
\newcommand{\Pdd}{\mbox{$\partial$ \hspace{-1.2 em} $/$}}
\newcommand{\Aslsh}{\slashed{A}}
\newcommand{\Bslsh}{\slashed{B}}
\newcommand{\Vslsh}{\slashed{V}}
\newcommand{\V}{\mathscr{V}}
\newcommand{\uslsh}{\slashed{u}}
\newcommand{\vslsh}{\slashed{v}}
\DeclareMathOperator{\Tr}{Tr}
\DeclareMathOperator{\tr}{tr}
\renewcommand{\O}{{\mathscr{O}}}
\renewcommand{\L}{{\mathcal{L}}}
\newcommand{\Sact}{{\mathcal{S}}}
\newcommand{\T}{{\mathcal{T}}}
\newcommand{\scrt}{{\mathfrak{t}}}
\newcommand{\M}{{\mathcal{M}}}
\newcommand\B{{\mathscr{B}}}
\newcommand{\U}{\text{\rm{U}}}
\newcommand{\SU}{\text{\rm{SU}}}
\newcommand{\g}{{\mathfrak{g}}}
\newcommand{\lf}{{\mathfrak{l}}}
\newcommand{\mf}{{\mathfrak{m}}}
\newcommand{\nf}{{\mathfrak{n}}}
\newcommand{\ef}{{\mathfrak{e}}}
\newcommand{\LR}{{L\!/\!R}}
\newcommand{\Cisc}{C^\infty_{\text{\rm{sc}}}}
\newcommand{\Cisco}{C^\infty_{\text{\rm{sc}},0}}
\newcommand{\Dir}{{\mathcal{D}}}
\DeclareMathOperator{\supp}{supp}
\renewcommand{\H}{\mathcal{H}}
\newcommand{\nuslsh}{\slashed{\nu}}
%\newcommand{\Dom}{{\mathscr{D}}}
\newcommand{\Lin}{\text{\rm{L}}}
\newcommand{\F}{{\mathscr{F}}}
\newcommand{\Tact}{{\mathcal{T}}}
\DeclareMathOperator{\Symm}{\mbox{\rm{Symm}}}
\newcommand{\nablaLC}{\nabla^\text{\tiny{\tt{LC}}}}
\newcommand{\DLC}{D^\text{\tiny{\tt{LC}}}}
\newcommand{\xis}{\slashed{\xi}}
\DeclareMathOperator{\norm}{| \hspace*{-0.1em}| }% \hspace*{-0.1em}|}
\newcommand{\K}{{\mathscr{K}}}
\newcommand{\frakD}{\mathfrak{D}}
\newcommand{\D}{\mathscr{D}}
\DeclareMathOperator{\re}{Re}
\DeclareMathOperator{\im}{Im}
\newcommand{\p}{\mathfrak{p}}
\newcommand{\SD}{\tilde{\Sig}^\text{\tiny{\rm{D}}}}
\newcommand{\I}{{\mathcal{I}}}
\DeclareMathOperator{\Pexp}{Pexp}
\DeclareMathOperator{\PE}{PE}
\newcommand{\FHF}{\mathscr{F}^\text{\tiny{HF}}}
\DeclareMathOperator{\sign}{sign}
\newcommand{\Sig}{\mathscr{S}}
\newcommand{\tildek}{\tilde{k}}
\newcommand{\scrL}{\mycal L}
\newcommand{\scrM}{\mycal M}
\newcommand{\scrN}{\mycal N}
\newcommand{\scrS}{\mycal S}
\newcommand{\w}{\item[{\raisebox{0.125em}{\tiny $\blacktriangleright$}}]}
\newcommand{\x}{{\textit{\myfont x}}}
\newcommand{\y}{{\textit{\myfont y}}}
\newcommand*{\myfont}{\fontfamily{ppl}\selectfont}		%ppl
\newcommand{\UMNS}{U_\text{\tiny{MNS}}}
\newcommand{\scrU}{{\mathscr{U}}}
\newcommand{\scrV}{{\mathscr{V}}}
\newcommand{\rel}{\text{\rm{rel}}}
\newcommand{\vol}{\text{\rm{\tiny{vol}}}}
\DeclareFontFamily{OT1}{rsfso}{}
\DeclareFontShape{OT1}{rsfso}{m}{n}{ <-7> rsfso5 <7-10> rsfso7 <10-> rsfso10}{}
\DeclareMathAlphabet{\mycal}{OT1}{rsfso}{m}{n}
\DeclareMathOperator{\vleck}{\mathcal{V}}
\newcommand{\cliff}{\!\cdot\!}
\DeclareMathOperator{\grad}{grad}
\newcommand{\Jdiff}{\mathfrak{J}^\text{\rm{\tiny{diff}}}}
\newcommand{\Jtest}{\mathfrak{J}^\text{\rm{\tiny{test}}}}
\newcommand{\Jlin}{\mathfrak{J}_\text{\rm{\tiny{lin}}}}
%\newcommand{\Jvary}{\mathfrak{J}^\text{\rm{\tiny{vary}}}}
\newcommand{\Gdiff}{\Gamma^\text{\rm{\tiny{diff}}}}
\newcommand{\Gtest}{\Gamma^\text{\rm{\tiny{test}}}}
\newcommand{\Ghom}{\Gamma^\text{\rm{\tiny{hom}}}}
\newcommand{\Ctest}{C^\text{\rm{\tiny{test}}}}
\renewcommand{\u}{\mathfrak{u}}
\renewcommand{\v}{\mathfrak{v}}
%\newcommand{\w}{\mathfrak{w}}
\newcommand{\J}{\mathfrak{J}}
\newcommand{\UL}{{\mathcal{Z}}}		% Symbol for the unit lattice in the lattice example
\newcommand{\itemD}{\item[{\raisebox{0.125em}{\tiny $\blacktriangleright$}}]}
\newcommand{\bei}{\begin{itemize}[itemsep=.2em]}
\newcommand{\eni}{\end{itemize}}
\DeclareMathOperator{\Pe}{Pe}

\newcommand\mpar[1]{\marginpar {\flushleft\sffamily\footnotesize #1}}
\setlength{\marginparwidth}{3.0cm}
\newcommand\Sean[1]{\marginpar {\flushleft\sffamily\footnotesize {\textcolor{Green}{Sean}}: #1}}
\newcommand\Johannes[1]{\marginpar {\flushleft\sffamily\footnotesize {\textcolor{CornflowerBlue}{Johannes}}: #1}}
\newcommand\bookmark{\marginpar {\flushleft\sffamily\footnotesize bookmark {\tt{Felix}}}}
% \definecolor{darkblue}{RGB}{0,91,163} %Farbe, um Änderungen zu markieren.


\newcommand{\Joh}[1]{\textcolor{CornflowerBlue}{#1}\mpar{\qquad$\circ$}} %Befehl um Änderungen zu markieren, Farbe für Johannes
\newcommand{\Sea}[1]{\textcolor{Green}{#1}\mpar{\qquad$\circ$}} %Befehl um Änderungen zu markieren, Farbe für Sean
\newcommand{\Skizze}{\textcolor{CornflowerBlue}{[Skizze]}}
\newcommand{\texit}[1]{\textit{#1}} % ;)
\newcommand{\itm}{\item[]}

% % Category theory 
% \newcommand{\cat}[1]{\mathbf{#1}}
% \newcommand{\Stoch}{\cat{Stoch}}
% \newcommand{\Set}{\cat{Set}}
\newcommand{\SProb}{\cat{SProb}}
% \newcommand{\id}[1]{\ensuremath{\mathrm{id}_{#1}}}

% % Discarding Symbols
% \newcommand{\discard}[1]{\ensuremath{\tinygroundnew_{#1}}}
% \newcommand{\discardflip}[1]{\ensuremath{\tinygroundflipnew_{#1}}}


%
% Tikzit Diagrams
%
\usepackage{tikz,xypic}
\usetikzlibrary{decorations.pathreplacing,decorations.markings,arrows.meta,backgrounds,shapes}
\usetikzlibrary{circuits.ee.IEC}
\pgfdeclarelayer{edgelayer}
\pgfdeclarelayer{nodelayer}
\pgfsetlayers{background,edgelayer,nodelayer,main}
\tikzstyle{none}=[inner sep=0mm]
\tikzstyle{every loop}=[]
\tikzstyle{mark coordinate}=[inner sep=0pt,outer sep=0pt,minimum size=3pt,fill=black,circle]

%Warning: this is also defined in the main preamble
\tikzset{arrow/.style={decoration={
    markings,
    mark=at position #1 with \arrow{>[length=2pt, width=3pt]}},
    postaction=decorate},
    reverse arrow/.style={decoration={
    markings,
    mark=at position #1 with {{\arrow{<[length=2pt, width=3pt]}}}},
    postaction=decorate}
}


% GROUND 
\tikzstyle{upground}=[circuit ee IEC,thick,ground,rotate=90,scale=1.5]

\newcommand{\minsym}{\mathsf{min}}
\newcommand{\sys}[1]{\ensuremath{\mathbf{#1}}}
% \newcommand{\sys}[1]{\ensuremath{\underline{#1}}}
\newcommand{\sysS}{\sys{S}}
\newcommand{\caus}{\mathsf{caus}}
\newcommand{\eff}{\mathsf{eff}}
\newcommand{\TPM}[2]{\ensuremath{T(#1,#2)}}
\newcommand{\TPMi}[2]{\ensuremath{T_i(#1,#2)}}
\newcommand{\TPMp}[2]{\ensuremath{p(#2 \mid #1)}}
\newcommand{\TZdir}{T^{\overset{\rightarrow}{Z}}}

\newcommand{\sysSZdir}{\sysS^{\overset{\rightarrow}{Z}}}
\newcommand{\crep}{\ensuremath{p_{\caus}}}
\newcommand{\exrep}{\ensuremath{\overline{R^{(x)}}}}
\newcommand{\eRerep}{\ensuremath{\overline{R^{\eff}}}}
\newcommand{\eRcrep}{\ensuremath{\overline{R^{\caus}}}}
\newcommand{\gxrep}{\ensuremath{R^{(x)}}}
\newcommand{\gerep}{\ensuremath{R^\eff}}
\newcommand{\gcrep}{\ensuremath{R^\caus}}
\newcommand{\uxrep}{\ensuremath{R^{(x)}_\mathsf{u}}}
\newcommand{\phicaus}{\ensuremath{\phi^{\caus}}}
\newcommand{\phieff}{\ensuremath{\phi^{\eff}}}
\newcommand{\erep}{\ensuremath{p_{\eff}}}
\newcommand{\uerep}{\ensuremath{{p_{\eff}}^u}} %Unconstrained effect repertoire
\newcommand{\eerep}{\ensuremath{\overline{p_{\eff}}}} %Extended effect repertoire
\newcommand{\ucrep}{\ensuremath{{p_{\caus}}^u}} %Unconstrained cause repertoire
\newcommand{\ecrep}{\ensuremath{\overline{p_{\caus}}}} %Extended cause repertoire

\newcommand{\crepp}[2]{\ensuremath{p_{\caus}}(#2 \mid #1)}
\newcommand{\erepp}[2]{\ensuremath{p_{\eff}}(#2 \mid #1)}
\newcommand{\ccaus}{P^\mathsf{c}}
\newcommand{\ceff}{P^\mathsf{e}}
\newcommand{\concept}[1]{\mathbb{C}(#1)}% {\mathsf{Concept}(#1)}
\newcommand{\con}[2]{\mathbb{C}_{#1}(#2)}
% {\mathsf{Concept}_{#1}(#2)}
\newcommand{\CEStruc}[2]{\mathbb{C}(#1,#2)}

% \newenvironment{picc}[1][]
% {\begin{aligned}\begin{tikzpicture}[font=\tiny,#1]}
% {\end{tikzpicture}\end{aligned}}


% \newcommand{\tinygroundnew}{
% \smash{
% % \raisebox{-2pt}
% {\hspace{-3pt}
% \ensuremath{
% \begin{picc}[scale=0.5] 
%     \node[upground, xscale=0.8, yscale=0.7] (1) at (0,0) {};
%     \draw (0,-0.2) to (0,-0.4);
% \end{picc}
% }\hspace{-1pt}}}}


% \newcommand{\tinygroundflipnew}{
% \smash{
% % \raisebox{1pt}
% {\hspace{-3pt}
% \ensuremath{
% \begin{picc}[xscale=0.5, yscale=-0.5] 
%     \node[upground, xscale=0.8, yscale=-0.7] (1) at (0,0.10) {};
%     \draw (0,-0.03) to (0,-0.31);
% \end{picc}
% }\hspace{-1pt}}}}

\newcommand{\beq}{\begin{equation}}
\newcommand{\eeq}{\end{equation}}
\newcommand{\Proof}{\begin{proof}}
\newcommand{\QED}{\end{proof} \noindent}
